\documentclass[14pt,a4paper]{extarticle}
\usepackage[T2A]{fontenc}
\usepackage[utf8]{inputenc}
\usepackage[english,russian]{babel}
\usepackage{lmodern}
\usepackage{geometry}
\usepackage{array}
\usepackage{tabularx}
\usepackage{longtable}
\usepackage{booktabs}
\usepackage{hyperref}
\usepackage{indentfirst}
\usepackage{microtype}
\usepackage{graphicx}
\usepackage{setspace}
\usepackage{listings}
\usepackage{graphicx}
\usepackage{caption}
\usepackage{float}
\usepackage{tocloft}
\usepackage{titlesec}
\usepackage{newtxtext,newtxmath}

\titleformat{\section}        % форматируем \section
  {\normalfont\bfseries}      % стиль: обычный + жирный
  {\thesection.}              % номер секции с точкой
  {2mm}                       % расстояние между номером и текстом
  {}                           % перед заголовком можно добавить что-то
\titlespacing*{\section}{1.5cm}{*3}{*1}  % {отступ слева}{отступ сверху}{отступ снизу}

\captionsetup{justification=centering}
\renewcommand{\ttdefault}{cmtt}
\setlength{\parindent}{3em}
\newgeometry{left=2.5cm,right=2.5cm,top=3cm,bottom=2.5cm}

\renewcommand{\contentsname}{\hspace*{4em}{\smallСодержание}}
\setlength{\cftbeforesecskip}{0.5em}
\renewcommand{\cftsecleader}{\cftdotfill{\cftdotsep}}
\renewcommand{\cftsecaftersnum}{.\hspace{2mm}}

\title{Отчёт по "Организация работы ИТ-службы"}
\author{Выполнил студент группы М-РИТ-25-1, Станиславчук Сергей}

\begin{document}

\setstretch{1.5}
\noindent
\begin{minipage}[t]{0.01\textwidth}
    \includegraphics{~/programmer/++Programmer/5_1/LSTU.jpg}
\end{minipage}%
\begin{minipage}[c]{1.1\linewidth}
    \centering
    {\small
    МИНИСТЕРСТВО НАУКИ И ВЫСШЕГО ОБРАЗОВАНИЯ РФ \\[1em]
    ФЕДЕРАЛЬНОЕ ГОСУДАРСТВЕННОЕ БЮДЖЕТНОЕ \\
    ОБРАЗОВАТЕЛЬНОЕ УЧРЕЖДЕНИЕ ВЫСШЕГО ОБРАЗОВАНИЯ \\
    «ЛИПЕЦКИЙ ГОСУДАРСТВЕННЫЙ ТЕХНИЧЕСКИЙ УНИВЕРСИТЕТ»
    }
\end{minipage}%

\vspace{3em}

\noindent
\hspace{1em}%
\begin{minipage}[t]{0.8\textwidth}
\textbf{Институт} \hspace{2em} компьютерных наук \\
\textbf{Кафедра} \hspace{2em} автоматизированных систем управления
\end{minipage}


\vspace{6em}

\begin{center}
\textbf{Практическая работа №5} \\
по дисциплине \\ <<Организация работы ИТ-службы>>
\end{center}

\vspace{6em}

\noindent
\begin{tabular*}{\textwidth}{@{\extracolsep{\fill}} l c r }
    Студент М-РИТ-25-1 & \makebox[4cm]{\hrulefill} & Станиславчук С.М. \\[2em]
Руководитель & \\
К.т.н., доцент & \makebox[4cm]{\hrulefill} & Ведищев В.В. \\
\end{tabular*}

\vfill
\begin{center}
    Липецк, 2026
\end{center}

\thispagestyle{empty}

\newpage
\tableofcontents
\newpage

\section{Эффективное планирование и управление потоком изменений.}

Любой IT-сервис постоянно изменяется. Для VDS-хостинга источниками изменений
являются обновление платформы виртуализации, модернизация систем хранения и
сети, рост числа заказчиков, новые тарифы и требования безопасности. Изменением
считается добавление, модификация или удаление элементов, способных повлиять на
услугу.

Цель управления изменениями -- обеспечить обоснованность,
планирование, согласование и приемлемый риск. Чрезмерный контроль создает
очередь и подталкивает сотрудников к обходу процедур, поэтому глубина
согласования должна соответствовать сложности работы.

Стандартные изменения выполняются по утвержденной процедуре: создание
виртуального сервера по шаблону, подключение диска, изменение расписания
резервного копирования. Нормальные требуют оценки влияния и согласования
коллегиальным органом по изменениям (CAB): обновление гипервизора, расширение
системы хранения, ввод тарифа. Экстренные выполняются при аварии или критической
уязвимости по ускоренной процедуре с последующим анализом.

Планирование начинается с единой очереди запросов (RFC), содержащих причину,
ожидаемый результат, затрагиваемые компоненты, срок, ответственного и оценку
риска. Поток регулируется приоритизацией по ценности и риску, ограничением
одновременно выполняемых работ, дроблением крупных изменений на проверяемые
этапы и стандартизацией повторяющихся операций.

Календарь изменений и окна обслуживания позволяют согласовывать работы между
эксплуатацией, сетью, безопасностью и поддержкой и уведомлять заказчиков. В
периоды пиковой нагрузки вводится заморозка изменений, при которой допускаются
только экстренные работы. Повторяющиеся операции автоматизируются через шаблоны
и описание конфигурации как кода.

Эффективность процесса оценивается по пропускной способности, времени выполнения
изменения, доле успешных изменений, возрасту очереди и соблюдению окон
обслуживания.

Пример: обновление гипервизора на кластере из нескольких десятков узлов.
Оформляется RFC, проверяется совместимость, обновление тестируется и согласуется
CAB. Затем узлы обновляются по одному: виртуальные машины переносятся, узел
выводится в обслуживание, обновляется и возвращается в кластер. Доступность и
производительность контролируются, предусмотрен откат.

\section{Управление рисками, связанными с изменением услуг.}

Риск -- функция, включающая в себя вероятность неблагоприятного события и его последствий,
В отличие от проблемы, риск только может реализоваться, поэтому к нему применяются предупреждающие меры.

Для каждого значимого изменения риски выявляются, оцениваются по вероятности и
влиянию, для них выбирается способ обработки, назначается владелец и организуется
контроль. Результаты фиксируются в реестре рисков или карточке изменения.

Источники риска: технические (несовместимость ПО, ошибки конфигурации, отказ
оборудования), операционные (нехватка персонала, ошибки выполнения),
бизнес-риски, информационная безопасность, правовые и финансовые факторы,
внешние поставщики -- дата-центр, операторы связи, поставщики оборудования.

Для качественной оценки используется матрица 4$\times$4: вероятность и влияние
оцениваются от 1 до 4, уровень риска определяется как $R = P \times I$. Оценки
1--3 считаются низкими, 4--8 средними, 9--16 высокими; высокий риск требует плана
обработки и согласования CAB. Пример матрицы приведен в таблице~\ref{tab:risk}.

\begin{table}[H]
    \centering
    \caption{Матрица оценки риска изменения}
    \label{tab:risk}
    \small
    \begin{tabular}{|l|c|c|c|c|}
        \hline
        \textbf{Вероятность / влияние} & \textbf{1} & \textbf{2} &
        \textbf{3} & \textbf{4} \\
        \hline
        Высокая (4) & 4 & 8 & 12 & 16 \\
        \hline
        Средняя (3) & 3 & 6 & 9 & 12 \\
        \hline
        Низкая (2) & 2 & 4 & 6 & 8 \\
        \hline
        Очень низкая (1) & 1 & 2 & 3 & 4 \\
        \hline
    \end{tabular}
\end{table}

Способы обработки: уклонение (отказ от изменения), снижение (тестирование,
резервирование, поэтапное развертывание), передача (договор, страхование) и
принятие остаточного риска с фиксацией решения.

План отката определяет условия признания изменения неуспешным, порядок возврата,
ответственных и допустимое время. Триггеры должны быть измеримыми: время отклика,
число ошибок, длительность недоступности, повреждение данных. До работ проверяется
восстановление резервных копий; человеческий фактор снижают чек-листы, разделение
обязанностей, принцип четырех глаз и ограничение прав.

Примеры рисков VDS: несовместимость гипервизора с гостевыми операционными
системами, падение производительности при перестроении массивов СХД, ошибки
маршрутизации при изменении сети. Мерами служат матрица совместимости, тестовый
кластер, работы в период низкой нагрузки, резервные копии и резервный канал
управления.

Пример: перенос профилей без проверки объема данных и состояния дисков привел
к росту времени входа, а замена сети не помогла из-за сбойного накопителя.
Аналогично миграция виртуальных машин без проверки хранилища может привести к
деградации сервиса. Реестр рисков пересматривается после изменений и инцидентов.

\section{Успешное внедрение изменений, отражающих новые требования бизнеса.}

Успешное изменение -- достижение бизнес-результата, который принимается
пользователями и устойчиво поддерживается. Рассмотрим внедрение услуги
резервного копирования виртуальных серверов: клиенты не имеют собственных
средств защиты данных, а часть заказчиков размещает критичные системы.

Требования: RPO (допустимая потеря данных) и RTO (время восстановления), сроки
хранения, поддерживаемые ОС, частота копирования, самостоятельное восстановление,
шифрование, разграничение доступа, оплата и ответственность сторон. Четкие
требования не позволяют создать решение, не соответствующее ожиданиям.

При проектировании определяются серверы резервного копирования, емкость СХД,
пропускная способность сети, интеграция с гипервизором, шифрование и контроль
доступа. Копии не должны храниться на той же СХД, что и основные данные.
Учитывается влияние копирования на производительность виртуальных машин.

Тестирование проверяет полные и инкрементальные копии, восстановление файлов и
целых виртуальных машин, прерывание копирования, нехватку места, отказ компонента
и безопасность. Пилот на ограниченной группе клиентов позволяет оценить
фактическую нагрузку, время копирования, удобство интерфейса и качество
поддержки.

Промышленное внедрение включает обновление каталога и биллинга, регистрацию в
Service Desk, мониторинг, документацию, обучение поддержки и инструкции для
клиентов. Критерии готовности: успешное тестирование, отсутствие критических
дефектов, готовность мониторинга и резервных мощностей, документация, обучение,
согласование условий и возможность отката. Если критический критерий не выполнен,
запуск переносится или ограничивается пилотом.

Критерии успеха: не менее 99\% успешных копий, 100\% успешных тестовых
восстановлений, RPO не более 24 часов (для критичных систем -- 1 час), RTO не
более 4 часов, влияние на производительность виртуальных машин не более 5\%,
подключение не менее 15\% клиентов за три месяца.

После запуска устанавливается период повышенного внимания, затем сравниваются
план и факт, анализируются инциденты и обратная связь. Типичные причины неудач --
нечеткие требования, отсутствие владельца, недооценка нагрузки, отсутствие пилота
и неготовность поддержки.

\section{Обеспечение качественных и актуальных знаний и информации по инфраструктуре и услугам.}

Знания являются таким же активом, как оборудование и программное обеспечение.
Без актуальной схемы сети или порядка восстановления специалист тратит больше
времени и чаще ошибается.

Система управления знаниями об услугах (SKMS) объединяет каталог услуг и SLA,
архитектурные схемы, описания конфигурационных единиц, регламенты, известные
ошибки, инструкции, договоры и отчеты по инцидентам. Для VDS выделяются
документация по услуге и инфраструктуре (сеть, IP-адресация, VLAN, кластеры,
СХД), регламенты (создание и миграция ВМ, вывод узла, замена диска,
восстановление), база поддержки, инструкции клиентов и разборы инцидентов.

Качество знания определяется точностью, полнотой, однозначностью, актуальностью,
доступностью, удобством поиска, краткостью и структурированностью. Текст
содержит цель, область применения, пошаговые действия, ожидаемый результат и
ограничения. Пароли, ключи и персональные данные в открытом виде не размещаются;
для конфиденциальных сведений устанавливаются права доступа.

Жизненный цикл включает создание, проверку, публикацию, использование,
обновление и архивирование. У статьи есть владелец, версия, дата, область
применения и связи с конфигурационными единицами. Регламенты критичных операций
пересматриваются не реже раза в квартал и после значимых изменений, инструкции
клиентов -- при изменении тарифов, схемы -- после изменения инфраструктуры.

Обновление документации -- обязательное условие закрытия изменения. Инциденты и
проблемы пополняют базу знаний: повторяющееся решение становится статьей,
серьезный инцидент -- разбором с хронологией, причинами и мерами. Полезно
фиксировать знания непосредственно в ходе обработки обращения.

Метрики: доля обращений, решенных с использованием базы знаний, обратная связь,
количество устаревших статей, время поиска и решения, повторные обращения, время
адаптации. Количество статей -- не показатель: неактуальные материалы ухудшают
поиск.

Пример: для услуги резервного копирования создаются инструкция клиента,
регламент, диагностическое дерево и описание ошибок. После пилота инструкция
дополняется примерами, после инцидента -- описанием причин. Без обновления
поддержка начнет работать по устаревшим правилам.

\section{Контроль за используемыми сервисными активами.}

Сервисные активы -- ресурсы и возможности организации, используемые для
предоставления услуг: оборудование, программное обеспечение, лицензии, договоры,
информационные ресурсы и персонал. Управление активами и конфигурациями
обеспечивает идентификацию, учет, контроль, отчетность и проверку достоверности
сведений.

Конфигурационная единица (CI) -- компонент, участвующий в предоставлении услуги.
Актив отражает ценность и учет ресурса, CI -- его участие в услуге и связи с
другими компонентами. Сведения хранятся в базе данных конфигураций (CMDB),
которая показывает компоненты, их взаимосвязи, владельцев, статусы и
местоположение.

Активы VDS: физические серверы, накопители, СХД, коммутаторы, маршрутизаторы,
межсетевые экраны, ИБП и стойки; виртуальные машины, шаблоны, сети, VLAN,
IP-адреса, тома и снимки; лицензии на виртуализацию, ОС, резервное копирование,
мониторинг и биллинг; договоры с дата-центром, операторами связи и поставщиками,
сертификаты и доменные имена.

Атрибуты: идентификатор, наименование, модель, серийный и инвентарный номера,
местоположение, владелец, статус, версия ПО, гарантия, стоимость, для виртуальных
объектов -- привязка к узлу, кластеру, сети и СХД.

Жизненный цикл включает закупку, получение, регистрацию, ввод в эксплуатацию,
эксплуатацию, обслуживание, модернизацию, вывод и списание. Статус меняется от
<<Заказано>> до <<Исключено из конфигурации>>; единая последовательность статусов
обязательна для всех подразделений.

Контроль: единые правила именования и маркировки, автоматическое обнаружение,
интеграция с закупками и финансовым учетом, инвентаризация и сверка с CMDB,
верификация записей, ограничение прав, связь с изменениями и инцидентами,
контроль лицензий и корректный вывод из эксплуатации.

При описании процесса определяются степень охвата (какие активы включаются) и
глубина детализации (какие атрибуты и связи фиксируются). Начинать целесообразно
с серверов и СХД, постепенно добавляя сеть, виртуальные ресурсы и договоры.

Пример: новый сервер регистрируется как <<Заказано>>, после поставки получает
инвентарный номер и статус <<Получено>>, перед включением в кластер фиксируются
место в стойке, сеть и адреса управления, после ввода -- статус <<В эксплуатации>>
и связи с ВМ, мониторингом и резервным копированием. При выводе ВМ переносятся,
данные уничтожаются, лицензии освобождаются, актив получает статус <<Исключено из
конфигурации>>.

CMDB помогает оценить влияние: при отказе узла видно, какие ВМ, заказчики и
сервисы затронуты; перед обновлением коммутатора -- какие кластеры и СХД
используют его. Учет стоимости, амортизации и договоров позволяет планировать
бюджет и оценивать себестоимость услуги; неучтенные ресурсы создают <<теневые>>
расходы и риск нарушения лицензий.

Метрики: полнота и точность учета, число неучтенных активов и доля автоматически
обнаруженных, соблюдение лицензий, загрузка ресурсов, стоимость владения, время
поиска конфигурационной единицы и ее связей.

\end{document}
